\documentclass[../RFC-HDFG-2026-003.tex]{subfiles}

\begin{document}

\section{Filter Plugin Class Extension}
\label{sec:class-extension}

\subsection{Blob Callbacks}
\label{sec:blob-callbacks}

Three new, NULL-able callbacks are added to \texttt{H5Z\_class3\_t}
(RFC-HDFG-2026-001): \texttt{write\_blob}, \texttt{read\_blob}, and
\texttt{close\_blob}. When NULL, the library uses a default
implementation built on \texttt{H5HG} (the existing global-heap
mechanism already used for large variable-length attribute data):
\texttt{H5HG\_insert} for writing, \texttt{H5HG\_read} for reading, and
\texttt{free()} for closing. Plugins implement these callbacks only when
they need custom on-disk layout --- for example, alignment requirements
for GPU DMA transfers directly out of the file buffer.

The three callbacks are independent of \texttt{set\_config}/\texttt{get\_config}:
a filter's parameter string continues to carry small, human-readable
configuration, while the blob carries the large, opaque payload. A
filter that needs both derives any small values it needs from dataset
properties via the existing \texttt{set\_local} callback, or encodes
them as a short header inside the blob itself; the blob content is
entirely filter-defined bytes, so this requires no new library
mechanism.

\begin{minted}{c}
/* Fixed-size on-disk locator for a blob's global-heap object.
 * Opaque to plugin authors; passed back unchanged between write_blob
 * and read_blob. */
typedef struct H5Z_blob_loc_t {
    haddr_t addr;   /* global heap collection address */
    size_t  idx;    /* object index within the collection */
} H5Z_blob_loc_t;

/* Called at H5Dcreate time to persist a blob alongside the pipeline
 * message.  If NULL, the library uses its default H5HG-based writer.
 * loc_out identifies the stored object for the later read_blob call. */
typedef herr_t (*H5Z_write_blob_func_t)(hid_t            file_id,
                                        const void       *buf,
                                        size_t            size,
                                        H5Z_blob_loc_t   *loc_out);

/* Called at H5Dopen time to recover a blob written by write_blob.
 * If NULL, the library uses its default H5HG-based reader.
 * *buf_out is allocated by the callback (or the library, in the
 * default case); the caller releases it via close_blob. */
typedef herr_t (*H5Z_read_blob_func_t)(hid_t            file_id,
                                       H5Z_blob_loc_t    loc,
                                       void            **buf_out,
                                       size_t           *size_out);

/* Called at H5Dclose time to release the in-memory blob buffer
 * returned by read_blob.  If NULL, the library calls free(). */
typedef herr_t (*H5Z_close_blob_func_t)(void *buf, size_t size);
\end{minted}

Added to \texttt{H5Z\_class3\_t}:

\begin{minted}{c}
    H5Z_write_blob_func_t       write_blob;  /* NULL-able */
    H5Z_read_blob_func_t        read_blob;   /* NULL-able */
    H5Z_close_blob_func_t       close_blob;  /* NULL-able */
\end{minted}

\texttt{H5Z\_blob\_loc\_t} is an opaque locator handed unchanged from
\texttt{write\_blob} to the pipeline-message encoder and later from the
decoder to \texttt{read\_blob}; plugin authors do not construct or
interpret its fields.

Add \texttt{has\_blob\_callbacks} (a \texttt{bool}) to
\texttt{H5Z\_class\_info\_t} (RFC-HDFG-2026-001) so
\texttt{H5Zget\_filter\_class\_info} can report whether a registered
filter implements the blob callbacks.

\subsection{New Internal State: \texttt{H5Z\_filter\_info\_t} Additions}
\label{sec:blob-internal-state}

\begin{minted}{c}
/* H5Zprivate.h -- internal only; not part of the public API */
struct H5Z_filter_info_t {
    /* ... existing fields ... */
    void            *aux_data;  /* in-memory blob; NULL if this filter has none */
    size_t           aux_size;  /* byte length of aux_data                     */
    H5Z_blob_loc_t   aux_loc;   /* on-disk locator; undefined until written    */
};
\end{minted}


\section{On-Disk Format}
\label{sec:blob-format}

The current \texttt{H5O\_PLINE\_VERSION\_2} per-filter layout:

\begin{center}
\texttt{[id:2][flags:2][cd\_nelmts:2][cd\_values:cd\_nelmts$\times$4][name:variable]}
\end{center}

This RFC's format adds a \texttt{has\_aux} bit after the name and, when
set, the fixed-size \texttt{H5Z\_blob\_loc\_t} encoding (address plus
heap-object index, the same fields \texttt{H5HG\_t} already carries
internally for other large-object references):

\begin{center}
\texttt{[id:2][flags:2][cd\_nelmts:2][cd\_values:cd\_nelmts$\times$4][name:variable]}\\
\texttt{[has\_aux:1][aux\_loc:variable if has\_aux]}
\end{center}

\textbf{Version number.} This RFC needs a pipeline-message format
bump. RFC-HDFG-2026-001 introduces no on-disk format change, so this RFC
is free to claim \texttt{H5O\_PLINE\_VERSION\_3} --- \emph{unless} some
other RFC has claimed that number first by the time this RFC is
finalized. Whoever finalizes this RFC should check for version-number
conflicts with any other pipeline-message format work in flight at that
time and claim the next free version.

The blob bytes themselves are stored as a global-heap (\texttt{H5HG})
object; the pipeline message records only the locator, not the bytes
inline. \texttt{H5O\_pline\_ver\_bounds[]} maps \texttt{H5F\_libver\_t}
to the pipeline message version; a DCPL with a blob-bearing filter
requires a libver upper bound that admits this RFC's pipeline-message
version, exactly as any other version-gated format feature does.

Runtime I/O is unaffected: \texttt{H5Z\_pipeline}, the hot path for
chunk I/O, touches none of this. The blob is configuration data, loaded
once at open time; compression/decompression performance is unchanged.


\section{Public API}
\label{sec:public-api}

\subsection{\texttt{H5Pappend\_filter\_blob}}
\label{sec:append-filter-blob}

\begin{minted}{c}
herr_t H5Pappend_filter_blob(
    hid_t         plist_id,
    H5Z_filter_t  id,
    unsigned      flags,
    const void   *buf,
    size_t        size
);
\end{minted}

\verysubsection{Description}
Appends filter \texttt{id} to the pipeline on \texttt{plist\_id}, exactly
as RFC-HDFG-2026-001's \texttt{H5Pappend\_filter} does, and attaches
\texttt{buf}/\texttt{size} as the filter's blob. The bytes are opaque to
the library; they are handed to the filter's \texttt{write\_blob}
callback (or the default \texttt{H5HG} writer) at \texttt{H5Dcreate}
time and recovered via \texttt{read\_blob} (or the default reader) at
\texttt{H5Dopen} time. \texttt{set\_config} is not invoked for this entry
point --- the blob path bypasses the parameter-string layer entirely,
since the whole point of the blob channel is to carry data that does
not fit through it.

\textbf{Decision: single entry point, no \texttt{H5Z\_params\_t}
variant.} An earlier draft considered adding \texttt{H5Z\_PARAMS\_BLOB}
as a third \texttt{H5Z\_params\_type\_t} tag so that blob configuration
could flow through \texttt{H5Pappend\_filter}'s tagged union. This RFC
does not do so: \texttt{H5Z\_params\_t} is designed around values that
fit comfortably in a compound literal, and forcing a multi-megabyte
buffer through that union would either double-store the buffer or
complicate its lifetime rules for no benefit. \texttt{H5Pappend\_filter\_blob}
is a dedicated function with a signature suited to its purpose;
\texttt{H5Z\_params\_type\_t} is unchanged by this RFC.

\verysubsection{Parameters}
\begin{itemize}
\item \texttt{plist\_id}---A dataset creation property list identifier.
\item \texttt{id}---The filter identifier. The registered filter must
  implement \texttt{write\_blob}/\texttt{read\_blob} or accept the
  library's default \texttt{H5HG}-based implementation.
\item \texttt{flags}---Same semantics as \texttt{H5Pappend\_filter}.
\item \texttt{buf}---Pointer to the blob bytes. The library copies these
  bytes into DCPL-owned storage before returning (\SectionRef{sec:blob-encode-decode});
  the caller's buffer may be freed or reused immediately after the call.
\item \texttt{size}---Length of \texttt{buf} in bytes. There is no fixed
  upper bound analogous to \texttt{H5Z\_CONFIG\_STRING\_MAX}; the
  practical limit is available memory and, in parallel jobs, the cost of
  the broadcast described in \SectionRef{sec:blob-parallel}.
\end{itemize}

\verysubsection{Returns}
Non-negative on success; negative on failure.

\verysubsection{Errors}
\begin{itemize}
\item \texttt{H5E\_NOFILTER}---The filter plugin could not be found or
  loaded.
\item \texttt{H5E\_BADVALUE}---\texttt{buf} is \texttt{NULL} with nonzero
  \texttt{size}, or \texttt{size} is zero with non-NULL \texttt{buf}.
\end{itemize}


\section{Dataset Create / Open / Close Lifecycle}
\label{sec:blob-lifecycle}

\textbf{Create} (\texttt{H5D\_\_create}), after \texttt{H5Z\_set\_local}
runs, for each filter with \texttt{aux\_data~!=~NULL}:
\begin{enumerate}
\item Call \texttt{filter->cls->write\_blob(file\_id, aux\_data, aux\_size,
  \&aux\_loc)}, or the default \texttt{H5HG} writer if
  \texttt{write\_blob} is \texttt{NULL}.
\item Store \texttt{aux\_loc} in \texttt{H5Z\_filter\_info\_t}.
\item Encode the pipeline message at this RFC's format version with
  \texttt{has\_aux} set and \texttt{aux\_loc} populated.
\end{enumerate}

\textbf{Open} (\texttt{H5D\_\_open}):
\begin{enumerate}
\item Decode the pipeline message; for each filter with \texttt{has\_aux},
  read \texttt{aux\_loc}.
\item Call \texttt{filter->cls->read\_blob(file\_id, aux\_loc, \&aux\_data,
  \&aux\_size)}, or the default \texttt{H5HG} reader if
  \texttt{read\_blob} is \texttt{NULL}.
\item Populate \texttt{aux\_data}/\texttt{aux\_size} in
  \texttt{H5Z\_filter\_info\_t}.
\end{enumerate}

\textbf{Close}: Call \texttt{close\_blob(aux\_data, aux\_size)}, or
\texttt{free()} if \texttt{close\_blob} is \texttt{NULL}, and clear the
fields.

\textbf{Delete} (dataset deleted, or the filter removed from a pipeline
being modified): the pipeline message's \texttt{delete} handler in
\texttt{H5Opline.c} calls \texttt{H5HG\_remove} on \texttt{aux\_loc} to
reclaim the heap space, analogous to how the fill-value message's
\texttt{delete} handler frees vlen data today (\texttt{H5Ofill.c}).


\section{Resolved Design Decisions}
\label{sec:decisions}

The four open design questions originally raised for this feature are
resolved as follows.

\subsection{Parallel I/O Write Protocol}
\label{sec:blob-parallel}

\noindent\textbf{Decision:} Rank~0 writes the blob; its locator is
broadcast to all ranks before the pipeline message is encoded.

\texttt{H5HG\_insert} has no MPI awareness: it allocates file space and
dirties a metadata-cache entry directly, outside the collective-I/O
coordination that object-header message appends go through. If every
rank in a collective \texttt{H5D\_\_create} called \texttt{H5HG\_insert}
independently, each would create a separate heap object at a different
address, with no way to agree on a single locator for the pipeline
message. The protocol:

\begin{minted}{c}
H5Z_blob_loc_t loc;

if (H5F_HAS_FEATURE(file, H5FD_FEAT_HAS_MPI)) {
    int mpi_rank;
    MPI_Comm_rank(H5F_mpi_get_comm(file), &mpi_rank);
    if (mpi_rank == 0)
        write_blob_or_default(file, aux_data, aux_size, &loc);
    MPI_Bcast(&loc, sizeof(loc), MPI_BYTE, 0, comm);
} else {
    write_blob_or_default(file, aux_data, aux_size, &loc);
}
/* all ranks now share the same locator; proceed to encode the
 * pipeline message collectively */
\end{minted}

This is the same pattern already used elsewhere in HDF5 for
single-writer allocations under \texttt{H5FD\_mpio} (e.g., superblock
writes). Calling \texttt{H5Pappend\_filter\_blob} on a handle opened for
parallel I/O is legal and must be followed by a collective
\texttt{H5D\_\_create}; when the DCPL is built before the file is opened
(the common pattern), the broadcast happens inside \texttt{H5D\_\_create},
invisible to the caller.

The broadcast is cheap regardless of blob size: only the fixed-size
locator (a few dozen bytes) is broadcast, never the blob bytes
themselves, so the cost added to a collective \texttt{H5Dcreate} is
independent of blob size.

\subsection{\texttt{H5Pencode}/\texttt{H5Pdecode} Policy}
\label{sec:blob-encode-decode}

\noindent\textbf{Decision:} \texttt{H5Pencode} serializes the blob's
bytes inline; \texttt{H5Pdecode} leaves the on-disk locator undefined
until the next \texttt{H5Dcreate}.

\texttt{H5P\_\_ocrt\_pipeline\_enc} encodes filter ID, flags, name, and
\texttt{cd\_values} --- all in-memory data; no file address appears in
an encoded DCPL. Fill-value encoding follows the same rule: a serialized
property must be self-contained, able to survive a buffer, a process
boundary, or a decode without the original file. \texttt{aux\_loc} is a
derived, file-specific value; \texttt{aux\_data} is the canonical
in-memory representation. \texttt{H5P\_\_ocrt\_pipeline\_enc} is
extended to append, after each filter's \texttt{cd\_values}:

\begin{center}
\texttt{[has\_aux:1][aux\_size:8 if has\_aux][aux\_data:aux\_size bytes if has\_aux]}
\end{center}

\texttt{H5P\_\_ocrt\_pipeline\_dec} reads those fields back into
\texttt{aux\_data}/\texttt{aux\_size} and leaves \texttt{aux\_loc}
undefined; a subsequent \texttt{H5D\_\_create} using the decoded DCPL
triggers \texttt{write\_blob} and assigns a real locator in whichever
file it is used with.

Serializing the locator instead of the bytes was rejected: it would make
an encoded DCPL valid only in its originating file, breaking every
\texttt{H5Pdecode}-then-\texttt{H5Dcreate} workflow in a different file.
Returning an error for blob-bearing DCPLs was also rejected as
unnecessarily restrictive, since \texttt{aux\_data} is always available
in memory after \texttt{H5D\_\_open} or after
\texttt{H5Pappend\_filter\_blob}. The tradeoff accepted is that encoded
DCPLs carrying a multi-megabyte blob are correspondingly large; this is
acceptable since \texttt{H5Pencode} is not a hot path and callers
invoking it on such a DCPL already know they hold unusual data.

\subsection{\texttt{h5repack} and \texttt{H5Pcopy} Behavior}
\label{sec:blob-repack}

\noindent\textbf{Decision:} No \texttt{h5repack}-specific code is
needed; \texttt{H5Pcopy} is extended to deep-copy \texttt{aux\_data}.

\texttt{h5repack} copies a normal (non-external) dataset's DCPL via
\texttt{H5Pcopy(dcpl\_in)} and does not call \texttt{H5HG} functions
directly. The existing lifecycle already makes this work automatically,
\emph{provided} \texttt{H5Pcopy} performs a deep copy of
\texttt{aux\_data}:

\begin{enumerate}
\item \texttt{H5Dopen} on the source populates \texttt{aux\_data} in the
  source DCPL's filter entry (via \texttt{read\_blob} or the default
  reader).
\item \texttt{H5Pcopy(dcpl\_in)} must \texttt{malloc} + copy the
  \texttt{aux\_data} buffer for each filter that has one, leaving
  \texttt{aux\_loc} undefined in the copy (the destination address is
  not yet known).
\item \texttt{H5Dcreate} with the copied DCPL in the destination file
  triggers \texttt{write\_blob} (or the default writer), producing a new
  heap object at a new locator in the destination file. Only the bytes
  travel; the source locator is never copied.
\end{enumerate}

The pipeline message's \texttt{copy} callback in \texttt{H5Opline.c} is
extended accordingly. When a filter is removed from a pipeline (e.g.,
\texttt{h5repack}'s \texttt{apply\_filters()} replacing a filter), the
pipeline message's destructor frees \texttt{aux\_data} the same way it
already must free other per-filter allocations.

\subsection{Storage Format: Naked \texttt{H5HG}}
\label{sec:blob-storage-format}

\noindent\textbf{Decision:} Blob bytes are stored as a plain
\texttt{H5HG} global-heap object, not a new HDF5 object message type.

A new message type (\texttt{H5O\_MSG\_FILTER\_BLOB}) would let
\texttt{h5dump} label the object (e.g., ``4~MB of SZ4 JIT source''), but
the payload remains filter-plugin-specific binary data that no generic
tool can render meaningfully --- the pipeline message already identifies
the filter, which is the useful part of that label. A new message type
also carries an ongoing maintenance cost (roughly a dozen callbacks:
encode, decode, copy, \texttt{raw\_size}, reset, free, delete, link, and
the pre/copy/post-copy-file hooks) for marginal tooling value. Existing
HDF5 precedent for ``large binary thing referenced from the object
header, recovered on demand'' is exactly \texttt{H5HG}, already used for
large variable-length attribute data. If a genuine tooling need for a
self-describing message type emerges after deployment, one can be
introduced in a later file-format version without breaking files
written under this design --- the pipeline message's \texttt{aux\_loc}
field would simply become one of several supported locator kinds.

\textbf{\texttt{del} callback obligation:} When a pipeline message is
deleted (dataset deleted, filter removed), \texttt{H5HG\_remove(file, \&loc)}
must be called to reclaim the heap space. This is implemented in the
pipeline message \texttt{delete} handler in \texttt{H5Opline.c},
analogous to how the fill-value message's \texttt{delete} handler frees
vlen data (\texttt{H5Ofill.c}).

\subsection{Use Case B: Reference to Another Dataset}
\label{sec:blob-usecaseb}

Use Case B (ROIBIN-SZ) stores an HDF5 reference to a user-visible mask
dataset as the filter's blob --- a small, fixed-size byte sequence
rather than a multi-megabyte payload, but otherwise an ordinary use of
\texttt{H5Pappend\_filter\_blob}. Dereferencing that reference to fetch
the mask's actual contents requires a file handle, which an earlier
draft proposed solving with a new \texttt{H5Z\_get\_file\_id()} accessor
or an additional \texttt{open\_dataset} callback.

\noindent\textbf{Decision:} Neither is needed. \texttt{read\_blob}
(\SectionRef{sec:blob-callbacks}) already receives \texttt{file\_id} ---
required for Use Case A's default \texttt{H5HG}-based storage --- and
that same handle is sufficient for Use Case B. A filter that stores a
reference as its blob dereferences it inside its own \texttt{read\_blob}
implementation using the supplied \texttt{file\_id}, fetches the
referenced dataset's contents once at open time, and caches them in
filter-private state (the \texttt{aux\_data} buffer, or memory it
allocates and tracks itself) for use during subsequent filter-callback
invocations. No per-chunk file access and no additional public API are
required; the same callback signature already specified for Use Case A
covers both use cases.


\section{Under Discussion: Retrieving a Default-Stored Blob}
\label{sec:blob-getter}

\begin{underdiscussion}
No decision is recorded for this section; it exists to solicit review
feedback before an accessor is added to the implementation.
\end{underdiscussion}

Prototype experience surfaced a gap in the default-storage path: a blob
written through the default \texttt{H5HG} writer is currently
\emph{write-only}. The library recovers the bytes into the pipeline at
\texttt{H5Dopen}, and they propagate correctly through \texttt{H5Pcopy},
\texttt{H5Pencode}/\texttt{H5Pdecode}, and \texttt{h5repack} --- but no
public API exposes them. An application cannot read a blob back from a
dataset's DCPL, and a filter cannot consume its own blob unless it
implements \emph{both} \texttt{write\_blob} and \texttt{read\_blob}: a
class that implements only \texttt{read\_blob} cannot interpret the
default writer's locator, because \texttt{H5HG} is library-private. Both
use cases in \SectionRef{sec:motivation} therefore require custom
callbacks today, which makes the default storage path considerably less
useful than it appears.

\noindent\textbf{Option under discussion:} a public accessor,

\begin{minted}{c}
herr_t H5Pget_filter_blob(
    hid_t     plist_id,
    unsigned  idx,        /* pipeline index, as in H5Pget_filter2       */
    void     *buf,        /* NULL for a size query                      */
    size_t   *size        /* in: capacity of buf; out: blob byte count */
);
\end{minted}

\noindent following the same size-query convention as
RFC-HDFG-2026-001's \texttt{H5Pget\_filter\_params\_by\_idx}. This would
let applications and tools (e.g., an \texttt{h5dump~-p} annotation that
a blob is present and its size) retrieve default-stored blobs, and ---
notably --- would let a \emph{filter} consume its blob without
implementing any storage callbacks: the \texttt{set\_local} callback
already receives the dataset's \texttt{dcpl\_id} and could fetch and
cache the blob there. That combination (default storage +
\texttt{set\_local} retrieval) would cover Use Case A with no
filter-side storage code at all.

Alternatives considered but not currently favoured: keeping the status
quo (filters that consume blobs must implement both storage callbacks),
or passing the blob into the filter callback through a new context
parameter (rejected for widening the hot-path signature).

\end{document}
